%%%%%%%%%%%%%%%%%%%%%%%%%%%%%%%%%%%%%%%%%%%%%%%%%%%%%%%%%%%%%%%%%%%%%%%%%%%%%%
%  Partie 6 -- Produire en France : la technologie est simple, la relocalisation possible par étapes
%  (2 diapositives : 17 = matrice des trois phases ; 18 = avec qui produire : injection,
%  puces de Grenoble, acier français)
%
%  Hypothèses chiffrées (à confirmer par devis) :
%  - carte française = 22 € (composants 14,5 + circuit 1,5 + assemblage 6), soit un coût de kit de
%    35,30 € (phase 2) puis 35,80 € (phase 3) ;
%  - phase 3 : boîtier injecté = pièce 1,50 € + moule amorti 1,50 € (≈ 33 k€ pour 22 000 kits,
%    moule en deux pièces : coque et capot) = 3,00 €, contre 2,50 € pour le boîtier imprimé en 3D ;
%  - les prix de 75 et 79 € ne contiennent PAS le surcoût éventuel de puces ou d'antennes françaises
%    (non chiffré) ;
%  - un fouet en acier de 3 mm sur 8 cm pèse ≈ 5 g : 22 000 kits ≈ 110 kg d'acier par an.
%%%%%%%%%%%%%%%%%%%%%%%%%%%%%%%%%%%%%%%%%%%%%%%%%%%%%%%%%%%%%%%%%%%%%%%%%%%%%%
\section{Produire en France}

% cellule « fabriqué en France » (coche rouge) et cellule « importé » (rond gris) : \cellok{x}{y}{texte}
\newcommand{\cellok}[3]{\pic[insarouge,scale=1.0] at ({#1-1.0},#2) {icone coche};
                        \node[anchor=west,font=\scriptsize] at ({#1-0.76},#2) {#3};}
\newcommand{\cellimp}[3]{\draw[insagris,thick] ({#1-1.0},#2) circle (0.13);
                         \node[anchor=west,font=\scriptsize,text=insagris] at ({#1-0.76},#2) {#3};}

%---------------------------------------------------------------------------
% Diapositive 17 : trois phases de relocalisation
%---------------------------------------------------------------------------
\begin{frame}{Technologie simple~: produisons en France}
  \centering
  \begin{tikzpicture}[font=\footnotesize]
    % en-têtes : trois phases, avec le prix public visé
    \node[pastille,draw=insarouge,fill=insarose,text=black,minimum width=2.8cm,minimum height=0.9cm] at (5.2,0)
      {An 1\\{\scriptsize\mdseries assemblé en France}\\{\scriptsize\color{insarouge}69\,€}};
    \node[pastille,draw=insarouge,fill=insarose,text=black,minimum width=2.8cm,minimum height=0.9cm] at (8.2,0)
      {Ans 2 et 3\\{\scriptsize\mdseries carte conçue ici}\\{\scriptsize\color{insarouge}75\,€}};
    \node[pastille,draw=insarouge,fill=insarose,text=black,minimum width=2.8cm,minimum height=0.9cm] at (11.2,0)
      {An 4 et après\\{\scriptsize\mdseries quasi tout ici}\\{\scriptsize\color{insarouge}79\,€}};
    % bandes de lecture (une ligne sur deux)
    \fill[insabeige] (0,-1.25) rectangle (12.7,-0.85);
    \fill[insabeige] (0,-2.05) rectangle (12.7,-1.65);
    \fill[insabeige] (0,-2.85) rectangle (12.7,-2.45);
    % libellés des lignes
    \node[anchor=west,font=\footnotesize\bfseries] at (0.1,-1.05) {Puce radio LoRa};
    \node[anchor=west,font=\footnotesize\bfseries] at (0.1,-1.45) {Processeur};
    \node[anchor=west,font=\footnotesize\bfseries] at (0.1,-1.85) {Carte électronique};
    \node[anchor=west,font=\footnotesize\bfseries] at (0.1,-2.25) {Boîtier};
    \node[anchor=west,font=\footnotesize\bfseries] at (0.1,-2.65) {Antenne};
    \node[anchor=west,font=\footnotesize\bfseries] at (0.1,-3.05) {Assemblage, flash, test};
    % puce radio : importée, puis fabriquée à Grenoble (objectif)
    \cellimp{5.2}{-1.05}{importée}\cellimp{8.2}{-1.05}{importée}\cellok{11.2}{-1.05}{Grenoble (visé)}
    % processeur : importé dans les trois phases
    \cellimp{5.2}{-1.45}{importé}\cellimp{8.2}{-1.45}{importé}\cellimp{11.2}{-1.45}{importé}
    % carte
    \cellimp{5.2}{-1.85}{module Heltec}\cellok{8.2}{-1.85}{conçue ici}\cellok{11.2}{-1.85}{conçue ici}
    % boîtier
    \cellok{5.2}{-2.25}{impression 3D}\cellok{8.2}{-2.25}{impression 3D}\cellok{11.2}{-2.25}{injection}
    % antenne
    \cellimp{5.2}{-2.65}{d'origine Heltec}\cellok{8.2}{-2.65}{acier, à valider}\cellok{11.2}{-2.65}{acier français}
    % assemblage
    \cellok{5.2}{-3.05}{France}\cellok{8.2}{-3.05}{France}\cellok{11.2}{-3.05}{France}
  \end{tikzpicture}

  \vspace{-0.45cm}
  \begin{columns}[T]
    \begin{column}{0.325\textwidth}
      \begin{carte}[height=1.05cm,valign=top]
        \raggedright\gros{+ 5 à 6\,€ par kit}\\[1pt]
        {\footnotesize 35,30\,€ au lieu de 30,30\,€.}
      \end{carte}
    \end{column}
    \begin{column}{0.325\textwidth}
      \begin{carte}[height=1.05cm,valign=top]
        \raggedright\gros{≈ 2 emplois}\\[1pt]
        {\footnotesize d'assemblage à 22\,000 kits.}
      \end{carte}
    \end{column}
    \begin{column}{0.325\textwidth}
      \begin{carte}[height=1.05cm,valign=top]
        \raggedright\gros{Né à Grenoble}\\[1pt]
        {\footnotesize LoRa~: inventé par Cycléo.}
      \end{carte}
    \end{column}
  \end{columns}
  \vspace{0cm}
  \begin{pointcle}\centering\footnotesize
    Carte de 5\,cm sur 2,5\,cm, composants du commerce~: rien d'irréalisable chez nous.
  \end{pointcle}
  \srcnote{Estimations à confirmer par devis~; prix de 75 et 79\,€ calculés pour garder une marge de 20 à 22\,€ par kit, hors surcoût éventuel de puces et d'antennes françaises. Sources~: \cite{actutem_lora,heltec_v3}}
\end{frame}

%---------------------------------------------------------------------------
% Diapositive 18 : avec qui produire ? injection plastique, puces de Grenoble, acier français
%---------------------------------------------------------------------------
\begin{frame}{Boîtier, puces, antennes~: avec qui produire~?}
  \begin{columns}[T]
    \begin{column}{0.6\textwidth}
      % schéma d'une presse à injecter : trémie, vis chauffante, moule en acier, pièce
      \begin{tikzpicture}[font=\footnotesize,scale=0.85]
        % trémie remplie de granulés
        \draw[insagris,thick,fill=insagris!10] (0.1,1.1) -- (1.3,1.1) -- (0.95,0.5) -- (0.45,0.5) -- cycle;
        \foreach \x/\y in {0.4/0.95,0.65/1.0,0.9/0.93,0.55/0.74,0.8/0.7} {\fill[insaorange] (\x,\y) circle (0.045);}
        \node[anchor=west,inner sep=1pt] at (1.45,0.92) {granulés};
        % fourreau, vis et colliers chauffants (rouges)
        \draw[insagris,thick,fill=insagris!15,rounded corners=2pt] (0,0) rectangle (3.3,0.5);
        \foreach \i in {0,...,9} {\draw[insagris,thin] ({0.2+0.3*\i},0.07) -- ({0.35+0.3*\i},0.43);}
        \foreach \x in {1.45,2.1,2.75} {\fill[insarouge] (\x,-0.05) rectangle (\x+0.12,0.55);}
        \fill[insagris!60] (-0.4,0.12) rectangle (0,0.38);
        \node[inner sep=1pt] at (1.65,-0.3) {vis chauffante};
        % buse et moule en acier : deux plaques et une cavité (matière en rouge)
        \draw[insagris,thick,fill=insagris!15] (3.3,0.13) -- (3.62,0.2) -- (3.62,0.3) -- (3.3,0.37) -- cycle;
        \draw[insagris,thick,fill=insagris!40] (3.62,-0.3) rectangle (5.4,0.25);
        \draw[insagris,thick,fill=insagris!40] (3.62,0.25) rectangle (5.4,0.8);
        \fill[insarouge!65] (3.62,0.22) rectangle (3.95,0.28);
        \draw[insarouge,fill=insarouge!65,rounded corners=1.5pt] (3.95,0.08) rectangle (5.05,0.42);
        \node[inner sep=1pt] at (4.5,1.0) {moule en acier};
        % pièce démoulée : le boîtier
        \draw[fluxr] (5.55,0.25) -- (6.25,0.25);
        \draw[insarouge,thick,fill=insarose,rounded corners=2pt] (6.4,-0.02) rectangle (7.6,0.52);
        \draw[insarouge,thick] (6.4,0.36) -- (7.6,0.36);
        \fill[insarouge] (7.38,0.18) circle (0.04);
        \node[inner sep=1pt] at (7.0,0.8) {boîtier};
      \end{tikzpicture}

      \vspace{0cm}
      {\footnotesize\raggedright
       \textbf{Injection~:} granulés fondus par la vis, injectés sous pression dans un moule en acier,
       refroidis, puis éjectés. Partenaires~: les plasturgistes d'Oyonnax (Plastics Vallée).\par}
    \end{column}
    \begin{column}{0.38\textwidth}
      \vspace{-0.3cm}
      \begin{carte}[height=1.05cm,valign=top]
        \raggedright\gros{5 à 50\,k€}\\[1pt]
        {\footnotesize pour un moule en acier.}
      \end{carte}
      \vspace{-0.15cm}
      \begin{carte}[height=1.05cm,valign=top]
        \raggedright\gros{0,50 à 5\,€}\\[1pt]
        {\footnotesize la pièce, hors moule.}
      \end{carte}
    \end{column}
  \end{columns}
  \vspace{0.02cm}
  \begin{columns}[T]
    \begin{column}{0.48\textwidth}
      \begin{encadre}{Puces~: l'écosystème de Grenoble}
        \begin{puces}
          \item Crolles grave en 22 puis 18\,nm (IoT, 5G)~: une puce radio, pas un processeur de pointe.
          \item ST, GlobalFoundries, Soitec. Piste~: STM32WL de ST, radio LoRa + processeur, soutenue par la communauté Meshtastic.
          \item Limite~: 22\,000 puces par an, trop peu pour une usine~; passer par un fabricant.
        \end{puces}
      \end{encadre}
    \end{column}
    \begin{column}{0.5\textwidth}
      \begin{encadre}{Antennes~: l'acier français}
        \begin{puces}
          \item ArcelorMittal~: 608 postes supprimés (Dunkerque, Florange, Basse-Indre…), mais 1,3\,Md€ en four électrique (2029).
          \item Basse-Indre~: acier plat étamé, soudable~: de quoi faire réflecteurs et plans de masse.
          \item ≈ 110\,kg d'acier par an, via des transformateurs. Cuivre, alu~: à comparer.
        \end{puces}
      \end{encadre}
    \end{column}
  \end{columns}
  \srcnote{Hypothèses~: moule ≈ 33\,k€ (1,50\,€ par kit) et pièce à 1,50\,€, soit 3,00\,€ contre 2,50\,€ imprimé~; fouet en acier de 5\,g × 22\,000 kits ≈ 110\,kg. Sources~: \cite{injection_cout,injection_procede,plastics_vallee,st_crolles,lmi_crolles,soitec,stm32wl,meshtastic_community,am_plan,am_four,am_basseindre,ttn_antennes}}
\end{frame}
